\documentclass[10pt,twocolumn]{article}

% --- Packages ---
\usepackage[utf8]{inputenc}
\usepackage[T1]{fontenc}
\usepackage{lmodern}
\usepackage[margin=1in]{geometry}
\usepackage{amsmath,amssymb}
\usepackage{graphicx}
\usepackage{booktabs}
\usepackage{multirow}
\usepackage{hyperref}
\usepackage{xcolor}
\usepackage{enumitem}
\usepackage{caption}
\usepackage{subcaption}
\usepackage{array}
\usepackage{tabularx}
\usepackage{longtable}
\usepackage{float}
\usepackage{authblk}
\usepackage{url}
\usepackage{microtype}
\usepackage{natbib}
\usepackage{xspace}

% --- Custom commands ---
\newcommand{\framework}{SecLens-R\xspace}
\newcommand{\seclens}{SecLens\xspace}
\newcommand{\tbd}{\textcolor{gray}{\textsf{TBD}}}
\newcommand{\placeholder}{\textcolor{gray}{\textsf{---}}}

\hypersetup{
    colorlinks=true,
    linkcolor=blue!60!black,
    citecolor=blue!60!black,
    urlcolor=blue!60!black
}

% --- Title ---
\title{\textbf{Role-Specific Evaluation of Large Language Models for Security Vulnerability Detection: A 250-Dimension Benchmark Framework}}

\author[1]{Shubo Haldar\thanks{Corresponding author. Email: \texttt{shubo.haldar@example.com}}}
\author[1]{Kashinath Kadaba Shrish}
\author[1]{Thiyagarajan M}
\affil[1]{Independent Researchers}

\date{}

\begin{document}

\maketitle

% ============================================================
% ABSTRACT
% ============================================================
\begin{abstract}
Large language models (LLMs) are increasingly deployed for automated security vulnerability detection, yet existing benchmarks reduce model capability to a single aggregate score. This oversimplification fails to address a fundamental organizational reality: different stakeholders---CISOs, heads of AI, security researchers, engineering leaders, and autonomous AI agents---have fundamentally different evaluation needs when selecting or deploying these models. We introduce \framework, a 250-dimension benchmark framework that extends the \seclens vulnerability detection benchmark with five role-specific evaluation lenses. Each role defines 50 tailored dimensions spanning ten measurement categories---detection accuracy, vulnerability knowledge depth, localization precision, reasoning quality, operational efficiency, tool-use proficiency, robustness, risk profile, scalability, and autonomy readiness---with independently calibrated weight vectors that sum to a composite Decision Score between 0 and 100. Our framework enables the same underlying benchmark data to produce distinct, actionable evaluations for each stakeholder. We present the complete dimension catalog, the weight assignment methodology, and the composite scoring formulation. We describe our integration with the two-layer \seclens evaluation pipeline (code-in-prompt and tool-use) and outline planned experiments across frontier and open-source LLMs. This work contributes the first multi-stakeholder evaluation framework for LLM-based vulnerability detection and establishes a foundation for organizationally-aware model selection.
\end{abstract}

\textbf{Keywords:} large language models, security vulnerability detection, benchmark, role-specific evaluation, multi-stakeholder decision framework

% ============================================================
% 1. INTRODUCTION
% ============================================================
\section{Introduction}
\label{sec:introduction}

The application of large language models (LLMs) to security vulnerability detection has emerged as one of the most promising directions in automated software security~\citep{pearce2022asleep,tony2023llmseceval,fang2024llmagents}. Models such as GPT-4~\citep{achiam2023gpt4}, Claude~\citep{anthropic2024claude}, and Gemini~\citep{team2023gemini} have demonstrated the ability to identify, classify, and localize vulnerabilities across diverse programming languages and CWE categories. This capability has significant implications for organizations seeking to augment their security programs with AI-driven tools.

However, the benchmarks used to evaluate these models---including CyberSecEval~\citep{bhatt2023cyberseceval}, SecEval~\citep{li2023seceval}, and others---share a critical limitation: they produce a single aggregate score that collapses the rich, multi-dimensional performance profile of a model into one number. While such scores are useful for leaderboard rankings, they are insufficient for the nuanced decisions that organizations actually face when adopting LLM-based security tools.

Consider the divergent needs of different organizational stakeholders:

\begin{itemize}[noitemsep,topsep=2pt]
    \item A \textbf{Chief Information Security Officer (CISO)} asks: ``Can I trust this model in my security program?'' The CISO prioritizes low false-negative rates on critical vulnerabilities, audit trail completeness, and regulatory compliance readiness. A model that excels at injection detection but silently misses authentication bypasses is unacceptable.
    \item A \textbf{Chief AI Officer (CAIO)} asks: ``Which model unlocks new capabilities while balancing risk and cost?'' The CAIO cares about the capability frontier---tool-use uplift, autonomous completion rates, and return on investment at scale.
    \item A \textbf{Security Researcher} asks: ``How deep and reliable is this model's vulnerability reasoning?'' Researchers need fine-grained CWE taxonomy mastery, evidence chain completeness (source, sink, data flow), and performance on rare vulnerability classes.
    \item A \textbf{Head of Engineering} asks: ``Will this help or hurt my team's velocity and code quality?'' Engineering leaders optimize for low false-positive rates (developer toil), CI/CD latency, actionable findings, and cost per pull request.
    \item An \textbf{AI-as-Actor} evaluation asks: ``Can this model operate autonomously on security tasks without human supervision?'' This lens evaluates tool-use mastery, self-correction, and decision quality in fully autonomous operation.
\end{itemize}

A single benchmark score cannot serve all five perspectives simultaneously. A model ranked first overall might be the worst choice for a CISO if it has a high overconfident miss rate on critical CWEs, or the worst choice for engineering if it generates excessive false positives that erode developer trust.

\subsection{Contributions}

We make the following contributions:

\begin{enumerate}[noitemsep,topsep=2pt]
    \item \textbf{A 250-dimension evaluation framework} comprising 50 role-specific dimensions for each of five stakeholder roles, organized across ten measurement categories.
    \item \textbf{A weighted composite scoring methodology} where each role has an independently calibrated weight vector that sums to 100, producing a role-specific Decision Score.
    \item \textbf{Integration with \seclens}, a two-layer benchmark that evaluates LLMs on vulnerability detection using both code-in-prompt (Layer 1) and tool-use (Layer 2) paradigms.
    \item \textbf{A complete dimension catalog} with computational definitions, data sources, and normalization strategies for all 250 dimensions.
\end{enumerate}

The remainder of this paper is organized as follows. Section~\ref{sec:related} reviews related work. Section~\ref{sec:framework} presents the framework design. Section~\ref{sec:catalog} provides the full dimension catalog. Section~\ref{sec:methodology} describes the evaluation methodology. Section~\ref{sec:experiments} outlines planned experiments. Section~\ref{sec:analysis} discusses expected analyses. Section~\ref{sec:discussion} addresses limitations and future work. Section~\ref{sec:conclusion} concludes.

% ============================================================
% 2. RELATED WORK
% ============================================================
\section{Related Work}
\label{sec:related}

\subsection{LLM Security Benchmarks}

Several benchmarks have been proposed for evaluating LLM capabilities in security contexts. CyberSecEval~\citep{bhatt2023cyberseceval} from Meta evaluates both the tendency of LLMs to generate insecure code and their ability to assist in cyberattacks, covering multiple programming languages and risk categories. SecEval~\citep{li2023seceval} focuses on code generation security across 130 CWE categories. PurpleLlama~\citep{bhatt2024purplellama} extends CyberSecEval with additional risk dimensions including prompt injection resistance and code interpreter abuse.

LLMSecEval~\citep{tony2023llmseceval} provides a dataset of security-relevant code completions based on MITRE CWE scenarios. SecurityEval~\citep{siddiq2022securityeval} offers a smaller but focused dataset for evaluating LLM-generated code for vulnerabilities. RepoSim~\citep{zhang2024reposim} evaluates LLMs in repository-level vulnerability detection with realistic multi-file contexts.

While these benchmarks provide valuable evaluation capabilities, they share two limitations relevant to our work: (1) they produce single aggregate scores or per-category scores without stakeholder-specific interpretation, and (2) they do not address the organizational decision context in which evaluation results are consumed.

\subsection{Role-Based Evaluation in Software Engineering}

The notion that different stakeholders require different evaluation lenses is well-established in software engineering. ISO/IEC 25010~\citep{iso25010} defines software quality from multiple viewpoints including users, developers, and operators. The Goal-Question-Metric (GQM) paradigm~\citep{basili1994gqm} formalizes the idea that metrics should be derived from stakeholder goals. In testing, risk-based approaches~\citep{amland1999riskbased} weight test cases by organizational impact rather than code coverage alone.

In the machine learning evaluation literature, Model Cards~\citep{mitchell2019modelcards} and Datasheets for Datasets~\citep{gebru2021datasheets} advocate for multi-stakeholder transparency. However, these approaches focus on documentation rather than providing different quantitative evaluation scores for different roles. Holistic Evaluation of Language Models (HELM)~\citep{liang2022helm} evaluates LLMs across multiple scenarios and metrics, but does not aggregate by stakeholder role.

\subsection{Multi-Stakeholder Decision Frameworks}

Multi-criteria decision analysis (MCDA)~\citep{velasquez2013mcda} provides formal methods for aggregating multiple evaluation dimensions with role-specific weight vectors. The Analytic Hierarchy Process (AHP)~\citep{saaty1990ahp} is a widely used MCDA technique for technology selection decisions. Our weight assignment methodology draws on these foundations while adapting them to the specific context of LLM security evaluation.

% ============================================================
% 3. FRAMEWORK DESIGN
% ============================================================
\section{Framework Design}
\label{sec:framework}

\subsection{Stakeholder Roles}
\label{sec:roles}

We define five stakeholder roles, each representing a distinct decision context in which LLM vulnerability detection capabilities are evaluated. Table~\ref{tab:roles} summarizes the roles and their core decision questions.

\begin{table}[t]
\centering
\caption{Stakeholder roles and their core evaluation questions.}
\label{tab:roles}
\small
\begin{tabular}{@{}p{2.0cm}p{5.0cm}@{}}
\toprule
\textbf{Role} & \textbf{Core Decision Question} \\
\midrule
CISO & ``Can I trust this model in my security program?'' \\
\addlinespace[3pt]
CAIO / Head of AI & ``Which model unlocks new capabilities while balancing risk and cost?'' \\
\addlinespace[3pt]
Security Researcher & ``How deep and reliable is this model's vulnerability reasoning?'' \\
\addlinespace[3pt]
Head of Engineering & ``Will this help or hurt my team's velocity and code quality?'' \\
\addlinespace[3pt]
AI as Actor & ``Can this model operate autonomously on security tasks?'' \\
\bottomrule
\end{tabular}
\end{table}

\textbf{CISO.} The Chief Information Security Officer is responsible for the organization's security posture, regulatory compliance, and risk management. The CISO lens prioritizes detection reliability on critical vulnerability classes (injection, RCE, authentication bypass), minimization of overconfident false negatives, audit trail completeness, and consistency across programming languages and CWE categories. The CISO is willing to accept higher costs and lower throughput in exchange for comprehensive, trustworthy coverage. Key unique dimensions include \emph{Critical Vulnerability Miss Prevention}, \emph{Compliance Report Readiness}, \emph{Board-Reportable Score}, and \emph{Risk Exposure Index}.

\textbf{CAIO / Head of AI.} The Chief AI Officer evaluates models from a strategic capability perspective, seeking to identify which model investments unlock the greatest organizational value per dollar. The CAIO lens balances raw capability against cost efficiency, deployment risk, and scalability. Unique dimensions include \emph{New Capability Index}, \emph{ROI Index}, \emph{Strategic Value Score}, and \emph{Business Case Strength}. The CAIO places significant weight on tool-use uplift and autonomous completion as indicators of frontier capability.

\textbf{Security Researcher.} The security researcher requires deep, technically rigorous evaluation. This lens prioritizes CWE taxonomy mastery (exact match rates, rare CWE identification), evidence chain completeness (source, sink, data flow analysis), cross-language consistency, and vulnerability-class-specific performance (injection mastery, XSS mastery, crypto weakness detection). Unique dimensions include \emph{Rare CWE Identification}, \emph{Source/Sink/Flow Identification Rates}, \emph{Injection/XSS/Auth/Crypto Class Mastery}, and \emph{Research Utility Index}.

\textbf{Head of Engineering.} The engineering leader optimizes for developer experience and CI/CD integration. This lens heavily penalizes false positives (developer toil), values fast response times (30-second and 60-second SLA compliance), and requires actionable findings (verdict + CWE + location). Unique dimensions include \emph{Cost Per PR}, \emph{Developer Trust Score}, \emph{CI Gate Readiness}, \emph{Developer Productivity Index}, and \emph{Team Adoption Readiness}.

\textbf{AI as Actor.} This lens evaluates the model's fitness for fully autonomous operation---a security agent that runs without human oversight. It prioritizes autonomous completion rates, tool-use mastery (adoption, efficiency, productivity), self-correction capability, and resistance to overconfident errors. Unique dimensions include \emph{Tool Call Productivity}, \emph{Agent Loop Efficiency}, \emph{Human Replacement Index}, \emph{Autonomous Triage Quality}, and \emph{Autonomous Scalability}.

\subsection{Dimension Taxonomy}
\label{sec:taxonomy}

Each role defines exactly 50 dimensions, yielding 250 role-specific dimensions in total. The dimensions are organized into ten categories that span the full evaluation space. Table~\ref{tab:categories} presents the category taxonomy with the approximate distribution of dimensions per role.

\begin{table}[t]
\centering
\caption{Ten measurement categories and their coverage focus.}
\label{tab:categories}
\small
\begin{tabular}{@{}clp{3.5cm}@{}}
\toprule
\textbf{ID} & \textbf{Category} & \textbf{Coverage Focus} \\
\midrule
A & Detection Accuracy & TPR, TNR, FPR, FNR, MCC, F1 \\
B & Vulnerability Knowledge & CWE identification, coverage, rare CWEs \\
C & Localization Precision & File, line-range, IoU accuracy \\
D & Reasoning Quality & Evidence chains, coherence, explanation \\
E & Operational Efficiency & Cost, tokens, wall time, throughput \\
F & Tool-Use \& Navigation & Tool calls, turns, search efficiency \\
G & Robustness \& Consistency & Parse rates, error rates, stability \\
H & Risk Profile & Critical CWEs, severity weighting, blind spots \\
I & Scalability \& Integration & Throughput, cost at scale, L1/L2 delta \\
J & Autonomy Readiness & Zero-guidance, self-correction, autonomous ops \\
\bottomrule
\end{tabular}
\end{table}

While the ten categories provide a shared organizational structure, each role defines its own 50 dimensions tailored to its decision context. Some dimensions are shared across roles (e.g., MCC, F1 score), while others are role-specific (e.g., the CISO's \emph{Board-Reportable Score} or the engineer's \emph{Cost Per PR}). Importantly, even when two roles measure the same underlying metric, they assign different weights reflecting different priorities.

\subsection{Weight Assignment Methodology}
\label{sec:weights}

Each role has a weight vector $\mathbf{w}^{(r)} = (w_1^{(r)}, w_2^{(r)}, \ldots, w_{50}^{(r)})$ where $r \in \{\text{CISO}, \text{CAIO}, \text{Researcher}, \text{Engineer}, \text{AI-Actor}\}$. The weights satisfy:

\begin{equation}
\sum_{i=1}^{50} w_i^{(r)} = 100 \quad \forall \, r
\label{eq:weight_constraint}
\end{equation}

Weight assignment follows a structured expert elicitation process:

\begin{enumerate}[noitemsep,topsep=2pt]
    \item \textbf{Role profiling.} For each role, we enumerate the top-5 decision criteria based on published job descriptions, industry frameworks (NIST CSF, ISO 27001, OWASP), and practitioner interviews.
    \item \textbf{Category allocation.} The 100 weight points are first distributed across the ten categories based on the role's priorities. For example, the CISO allocates 23\% to Risk Profile and 21\% to Detection Accuracy, while the AI-as-Actor allocates 23\% to Autonomy Readiness and 17\% to Tool-Use.
    \item \textbf{Intra-category distribution.} Within each category, weights are distributed among the constituent dimensions based on their decision relevance.
    \item \textbf{Normalization check.} The final vector is verified to sum to exactly 100.
\end{enumerate}

Table~\ref{tab:category_weights} presents the category-level weight distribution across all five roles.

\begin{table*}[t]
\centering
\caption{Category-level weight distribution across roles (weights sum to 100 per role).}
\label{tab:category_weights}
\small
\begin{tabular}{@{}l r r r r r@{}}
\toprule
\textbf{Category} & \textbf{CISO} & \textbf{CAIO} & \textbf{Researcher} & \textbf{Head Eng.} & \textbf{AI Actor} \\
\midrule
A: Detection Accuracy        & 21.0 & 12.0 & 12.0 & 14.0 &  9.5 \\
B: Vulnerability Knowledge   & 10.0 &  5.0 & 20.0 &  3.5 &  9.0 \\
C: Localization Precision    &  4.5 &  2.5 & 12.0 & 13.0 &  8.0 \\
D: Reasoning Quality         &  9.0 & 11.0 & 15.0 &  5.5 & 11.5 \\
E: Operational Efficiency    &  7.0 & 10.0 &  3.5 & 13.0 &  5.5 \\
F: Tool-Use \& Navigation    &  3.5 & 14.0 &  8.5 &  8.5 & 17.0 \\
G: Robustness \& Consistency & 10.0 & 10.0 &  5.0 & 13.0 & 13.0 \\
H: Risk Profile              & 23.0 & 12.0 & 13.0 &  8.0 & 10.5 \\
I: Scalability \& Integration&  7.0 &  8.0 &  5.5 & 12.0 &  8.5 \\
J: Autonomy Readiness        & 10.0 & 18.0 & 11.0 &  6.0 & 23.0 \\
\midrule
\textbf{Total}               & \textbf{100.0} & \textbf{100.0} & \textbf{100.0} & \textbf{100.0} & \textbf{100.0} \\
\bottomrule
\end{tabular}
\end{table*}

The weight distribution reveals clear role-specific priorities. The CISO allocates the most weight to Risk Profile (23\%) and Detection Accuracy (21\%), reflecting the primacy of comprehensive, reliable threat detection. The CAIO prioritizes Autonomy Readiness (18\%) and Tool-Use (14\%), seeking models that push the capability frontier. The Security Researcher concentrates on Vulnerability Knowledge (20\%) and Reasoning Quality (15\%), valuing depth of understanding. The Head of Engineering distributes weight more evenly, with emphasis on Detection Accuracy (14\%), Operational Efficiency (13\%), Robustness (13\%), Localization (13\%), and Scalability (12\%). The AI-as-Actor lens heavily weights Autonomy Readiness (23\%) and Tool-Use (17\%).

\subsection{Composite Decision Score}
\label{sec:composite}

For a model $m$ evaluated on role $r$, let $s_i^{(m)}$ denote the normalized score (in $[0, 1]$) on dimension $i$. The composite Decision Score is:

\begin{equation}
D^{(r)}(m) = \sum_{i=1}^{50} w_i^{(r)} \cdot s_i^{(m)}
\label{eq:decision_score}
\end{equation}

Since all $w_i^{(r)}$ sum to 100 and all $s_i^{(m)} \in [0, 1]$, the Decision Score is naturally bounded in $[0, 100]$.

\textbf{Normalization.} Most dimensions are naturally in $[0, 1]$ (rates, proportions). For unbounded dimensions (e.g., cost per task, tokens per task, wall time), we apply min-max normalization against the evaluated model cohort, with optional inversion for metrics where lower is better:

\begin{equation}
s_i^{(m)} = 1 - \frac{v_i^{(m)} - \min_k v_i^{(k)}}{\max_k v_i^{(k)} - \min_k v_i^{(k)}}
\label{eq:normalization}
\end{equation}

where $v_i^{(m)}$ is the raw value of dimension $i$ for model $m$, and the min/max are taken over all models in the evaluation cohort.

% ============================================================
% 4. DIMENSION CATALOG
% ============================================================
\section{Dimension Catalog}
\label{sec:catalog}

This section presents the complete 250-dimension catalog. Each role defines 50 dimensions with unique identifiers, descriptive names, and computational definitions grounded in data available from the \seclens evaluation pipeline.

\subsection{CISO Dimensions}
\label{sec:ciso_dims}

The CISO lens comprises 50 dimensions organized to answer: ``Can I trust this model in my security program?'' The dimensions emphasize detection reliability on critical vulnerability classes, risk governance metrics, compliance readiness, and operational stability. Table~\ref{tab:ciso_dims} presents the complete catalog.

\begin{table*}[t]
\centering
\caption{CISO dimension catalog (50 dimensions). All scores normalized to $[0, 1]$ unless noted.}
\label{tab:ciso_dims}
\small
\begin{tabular}{@{}r l p{7.5cm}@{}}
\toprule
\textbf{\#} & \textbf{Dimension} & \textbf{Description} \\
\midrule
C1  & Overall Detection Reliability       & Verdict accuracy across all tasks; top-line trust metric \\
C2  & Critical Vuln Detection Rate        & Recall on injection, RCE, auth bypass, XSS, deserialization CWEs \\
C3  & Critical Vuln Miss Prevention       & Complement of miss rate on critical CWE classes \\
C4  & False Alarm Suppression             & True negative rate; measures organizational noise reduction \\
C5  & Balanced Detection Accuracy         & (TPR + TNR) / 2; balanced view of detection quality \\
C6  & Overall MCC Score                   & Matthews Correlation Coefficient normalized to $[0,1]$ \\
C7  & Audit Trail Completeness            & Fraction of responses with source, sink, and data flow evidence \\
C8  & CWE Classification Accuracy         & Correct CWE-ID among true positive detections \\
C9  & Zero-Day Readiness                  & Performance on rare CWE categories ($\leq 3$ tasks); novel vuln proxy \\
C10 & Severity-Weighted Recall            & Recall weighted by assumed CWE severity tier (critical $3\times$) \\
C11 & Worst-Category Floor                & Minimum accuracy across CWE categories; no blind spots \\
C12 & Language Parity                     & $1 - (\max - \min)$ accuracy gap across programming languages \\
C13 & Overconfident Miss Prevention       & $1 -$ FN rate among FULL-parsed (high-confidence) responses \\
C14 & Operational Error Rate              & Fraction of tasks completing without crash or error \\
C15 & Parse Reliability                   & Fraction with parseable output (FULL or PARTIAL) \\
C16 & Structured Output Compliance        & Fraction with fully schema-compliant JSON output \\
C17 & Location Accuracy for Remediation   & Location score among correct TP verdicts \\
C18 & Finding Actionability               & Fraction of TPs with CWE + location (fully actionable) \\
C19 & Negative Predictive Value           & Among ``not vulnerable'' predictions, fraction correct \\
C20 & Positive Predictive Value           & Precision; among ``vulnerable'' predictions, fraction correct \\
C21 & Cost Per Task                       & Average USD per task \\
C22 & Cost Per True Positive              & Dollars per correctly detected vulnerability \\
C23 & Budget Predictability               & $1 - CV$ of cost distribution; consistent spend \\
C24 & Annual Projected Cost               & Projected cost for 10K tasks/year (raw USD) \\
C25 & Compliance Report Readiness         & Fraction of findings with CWE + evidence + location \\
C26 & Reasoning Transparency              & Fraction of responses with reasoning field populated \\
C27 & Cross-Run Consistency               & Placeholder; requires multi-seed comparison \\
C28 & Regression Risk                     & Placeholder; requires model version comparison \\
C29 & Supply Chain CWE Coverage           & Performance on dependency/supply-chain CWE classes \\
C30 & Authentication Bypass Detection     & Recall on auth/authz/session CWE categories \\
C31 & Injection Detection Rate            & Recall on injection-class vulnerabilities \\
C32 & XSS Detection Rate                  & Recall on cross-site scripting vulnerabilities \\
C33 & Data Exposure Detection             & Recall on information disclosure vulnerabilities \\
C34 & Cryptographic Weakness Detection    & Recall on crypto/random/hash weakness CWEs \\
C35 & Mean Time to Detect                 & Normalized wall time (${<}10$s$=1.0$, ${>}120$s$=0.0$) \\
C36 & Scanning Throughput                 & Tasks per minute, capped at 30 TPM $= 1.0$ \\
C37 & Incident Response Utility           & Composite: $0.5{\times}$TPR + $0.2{\times}$speed + $0.3{\times}$actionability \\
C38 & Vendor Risk Assessment Utility      & $0.5{\times}$language parity + $0.5{\times}$CWE breadth \\
C39 & Regulatory Mapping Accuracy         & CWE accuracy; CWEs map to PCI-DSS, SOC 2, HIPAA \\
C40 & Triage Noise Ratio                  & PPV; fraction of alerts worth investigating \\
C41 & SLA Compliance Rate                 & Fraction completing within 60 seconds \\
C42 & Token Efficiency                    & Accuracy per 1K tokens consumed \\
C43 & Multi-Language Coverage             & Fraction of languages with ${>}50\%$ accuracy \\
C44 & Patch Verification Accuracy         & Accuracy on post-patch negative tasks \\
C45 & SAST FP Rejection                   & Accuracy on SAST false positive tasks \\
C46 & Board-Reportable Score              & Severity-weighted F1 score \\
C47 & Risk Exposure Index                 & $0.6{\times}$critical detection + $0.4{\times}$overconfidence prevention \\
C48 & Defense-in-Depth Contribution       & Value-add alongside SAST; rare CWE performance proxy \\
C49 & Production Stability                & $0.5{\times}$error rate + $0.5{\times}$parse reliability \\
C50 & Overall CISO Confidence             & $0.4{\times}$balanced accuracy + $0.35{\times}$risk index + $0.25{\times}$stability \\
\bottomrule
\end{tabular}
\end{table*}

\subsection{CAIO / Head of AI Dimensions}
\label{sec:caio_dims}

The CAIO lens comprises 50 dimensions focused on strategic capability assessment, cost-efficiency, and organizational value creation. Table~\ref{tab:caio_dims} presents the complete catalog.

\begin{table*}[t]
\centering
\caption{CAIO / Head of AI dimension catalog (50 dimensions).}
\label{tab:caio_dims}
\small
\begin{tabular}{@{}r l p{7.5cm}@{}}
\toprule
\textbf{\#} & \textbf{Dimension} & \textbf{Description} \\
\midrule
A1  & Capability Score                    & Overall verdict accuracy; baseline capability signal \\
A2  & Full Finding Rate                   & Fraction of TP tasks with perfect 3/3 score \\
A3  & MCC (Normalized)                    & Matthews Correlation Coefficient, normalized to $[0,1]$ \\
A4  & Recall                              & True positive rate \\
A5  & Precision                           & Positive predictive value \\
A6  & F1 Score                            & Harmonic mean of precision and recall \\
A7  & Tool-Use Effectiveness              & Among tool-using tasks, fraction with score ${>}0$ \\
A8  & Tool-Use Uplift                     & Accuracy delta: tool-using vs non-tool-using tasks \\
A9  & Autonomous Completion Rate          & Fraction without error or parse failure \\
A10 & Multi-Turn Success Rate             & Accuracy on tasks requiring multiple conversation turns \\
A11 & Navigation Efficiency               & Fraction of L2 tasks solved with ${\leq}5$ tool calls \\
A12 & Reasoning Depth                     & Fraction with evidence chain (source + sink) \\
A13 & Zero-Guidance Capability            & Accuracy under minimal prompt guidance \\
A14 & Cross-Category Breadth              & Fraction of CWE categories with ${>}50\%$ accuracy \\
A15 & Language Breadth                    & Fraction of languages with ${>}50\%$ accuracy \\
A16 & Cost Per Task                       & Average USD per task \\
A17 & Cost Per Correct Finding            & USD per true positive detected \\
A18 & MCC Per Dollar                      & MCC divided by total cost; quality per dollar \\
A19 & Tokens Per Task                     & Average total tokens consumed \\
A20 & Accuracy Per 1K Tokens              & Correct verdicts per 1K tokens consumed \\
A21 & Cost at Scale (1K)                  & Projected cost for 1,000 tasks \\
A22 & Cost at Scale (10K)                 & Projected cost for 10,000 tasks \\
A23 & Wall Time Per Task                  & Average seconds per task \\
A24 & Throughput                          & Tasks per minute \\
A25 & Risk-Adjusted Return                & Capability ${\times}$ reliability / cost \\
A26 & Parse Compliance                    & FULL parse rate; structured output for integration \\
A27 & Error Rate (inverted)               & $1 -$ error rate; reliability for automated pipelines \\
A28 & Graceful Degradation                & $1 - |\text{common\_acc} - \text{rare\_acc}|$ \\
A29 & Self-Correction Capability          & Multi-turn task accuracy; iterative refinement \\
A30 & False Positive Prevention           & True negative rate \\
A31 & Critical Miss Prevention            & Detection rate on high-severity CWE categories \\
A32 & Worst-Category Risk                 & Minimum accuracy across categories; downside risk \\
A33 & Language Consistency                 & $1 -$ language accuracy gap \\
A34 & Overconfident Miss Prevention       & $1 -$ FN rate among FULL-parsed responses \\
A35 & Patch Verification                  & Post-patch task accuracy \\
A36 & Workflow Automation Readiness       & $0.4{\times}$auto + $0.3{\times}$parse + $0.3{\times}$error rate \\
A37 & New Capability Index                & $0.4{\times}$tool uplift + $0.3{\times}$multi-turn + $0.3{\times}$rare CWE \\
A38 & Competitive Differentiation         & Full finding rate; measures peak capability \\
A39 & ROI Index                           & F1 / cost\_per\_task; return per dollar \\
A40 & Reasoning Quality                   & Fraction with reasoning AND correct verdict \\
A41 & Confidence Calibration              & FULL parse + correct verdict rate \\
A42 & Integration Readiness               & $0.4{\times}$parse + $0.3{\times}$error + $0.3{\times}$autonomous \\
A43 & CWE Knowledge Depth                 & CWE accuracy among TPs \\
A44 & Location Precision                  & Location accuracy among TPs \\
A45 & Strategic Value Score               & $0.3{\times}$cap + $0.2{\times}$eff + $0.2{\times}$rel + $0.3{\times}$new \\
A46 & SAST Augmentation Value             & SAST FP task accuracy \\
A47 & Deployment Risk (inverted)          & $1 - (\text{error} + \text{parse\_fail}) / 2$ \\
A48 & Scalability Confidence              & $1 - CV$ of cost distribution \\
A49 & Business Case Strength              & $0.4{\times}$F1 + $0.3{\times}$efficiency + $0.3{\times}$risk \\
A50 & Overall CAIO Score                  & $0.35{\times}$capability + $0.25{\times}$quality + $0.2{\times}$risk + $0.2{\times}$cost \\
\bottomrule
\end{tabular}
\end{table*}

\subsection{Security Researcher Dimensions}
\label{sec:researcher_dims}

The Security Researcher lens defines 50 dimensions emphasizing vulnerability reasoning depth, CWE taxonomy mastery, evidence quality, and per-language/per-class performance. Table~\ref{tab:researcher_dims} presents the complete catalog.

\begin{table*}[t]
\centering
\caption{Security Researcher dimension catalog (50 dimensions).}
\label{tab:researcher_dims}
\small
\begin{tabular}{@{}r l p{7.5cm}@{}}
\toprule
\textbf{\#} & \textbf{Dimension} & \textbf{Description} \\
\midrule
R1  & Vulnerability Detection Recall      & True positive rate \\
R2  & Vulnerability Detection Precision   & Positive predictive value \\
R3  & F1 Score                            & Harmonic mean of precision and recall \\
R4  & MCC (Normalized)                    & Matthews Correlation Coefficient, $[0,1]$ \\
R5  & CWE Exact Match Rate               & Exact CWE-ID match among correct TP verdicts \\
R6  & CWE Coverage Breadth                & Fraction of CWE categories with ${\geq}1$ correct CWE \\
R7  & Rare CWE Detection                  & Verdict accuracy on CWEs with ${\leq}3$ tasks \\
R8  & Rare CWE Identification             & CWE match accuracy on rare CWE categories \\
R9  & Cross-Language CWE Consistency      & $1 - \text{StdDev}$ of CWE accuracy across languages \\
R10 & CWE Confusion Quality               & $1 -$ fraction of wrong CWE predictions \\
R11 & Location Accuracy (IoU${>}$0.5)     & Location score among correct TP verdicts \\
R12 & File Identification Rate            & Fraction of TPs with location field populated \\
R13 & Mean IoU Estimate                   & Estimated mean IoU (location$=1 \rightarrow 0.75$) \\
R14 & Pinpoint Rate                       & Proxy: same as location accuracy \\
R15 & Over-Scope Prevention               & $1 -$ fraction of location attempts that are wrong \\
R16 & Evidence Completeness               & Fraction with source + sink + data flow chain \\
R17 & Source Identification Rate           & Fraction of TPs with source identified \\
R18 & Sink Identification Rate             & Fraction of TPs with sink identified \\
R19 & Data Flow Completeness              & Fraction of TPs with non-empty flow chain \\
R20 & Reasoning Presence                  & Fraction with reasoning field populated \\
R21 & Reasoning Length (median)           & Median reasoning length in characters \\
R22 & Reasoning + Correct Verdict         & Fraction with reasoning AND correct verdict \\
R23 & FP Reasoning Quality                & Among FPs, fraction with reasoning present \\
R24 & Injection Class Mastery             & CWE accuracy on injection-class vulnerabilities \\
R25 & XSS Class Mastery                   & CWE accuracy on XSS vulnerabilities \\
R26 & Auth Class Mastery                  & CWE accuracy on auth/access control vulnerabilities \\
R27 & Crypto Class Mastery                & CWE accuracy on cryptographic weakness CWEs \\
R28 & Deserialization Mastery             & Detection rate on deserialization vulnerabilities \\
R29 & Worst CWE Category                  & Minimum accuracy across CWE categories \\
R30 & Best CWE Category                   & Maximum accuracy across CWE categories \\
R31 & Category Variance (inv.)            & $1 - \text{StdDev}$ of accuracy across categories \\
R32 & Python Accuracy                     & Accuracy on Python tasks specifically \\
R33 & JavaScript Accuracy                 & Accuracy on JavaScript/TypeScript tasks \\
R34 & Java Accuracy                       & Accuracy on Java tasks \\
R35 & C/C++ Accuracy                      & Accuracy on C and C++ tasks \\
R36 & Go Accuracy                         & Accuracy on Go tasks \\
R37 & Language Gap (inverted)             & $1 -$ (best language $-$ worst language) accuracy \\
R38 & Post-Patch Discrimination           & Accuracy on post-patch (patched code) tasks \\
R39 & SAST FP Discrimination              & Accuracy on SAST false positive tasks \\
R40 & Negative Task Accuracy              & Overall specificity on negative tasks \\
R41 & Complete Finding Rate (3/3)         & Fraction of TPs with perfect 3/3 score \\
R42 & Tool-Use for Deep Analysis          & Multi-turn task accuracy \\
R43 & Novel Pattern Detection             & Rare CWE detection rate (same as R7) \\
R44 & Graceful Degradation                & $1 - |\text{common\_acc} - \text{rare\_acc}|$ \\
R45 & Parse Success Rate                  & FULL parse fraction \\
R46 & Error Rate (inverted)               & $1 -$ error rate \\
R47 & Reproducibility Proxy               & Placeholder (0.5); requires multi-seed runs \\
R48 & Prompt Robustness                   & Placeholder (0.5); requires multi-preset runs \\
R49 & Research Utility Index              & $0.4{\times}$CWE + $0.3{\times}$evidence + $0.3{\times}$rare CWE \\
R50 & Overall Researcher Score            & Composite: detection, CWE, evidence, breadth, location \\
\bottomrule
\end{tabular}
\end{table*}

\subsection{Head of Engineering Dimensions}
\label{sec:engineer_dims}

The Head of Engineering lens defines 50 dimensions optimized for developer velocity, CI/CD integration, and signal-to-noise ratio. Table~\ref{tab:engineer_dims} presents the complete catalog.

\begin{table*}[t]
\centering
\caption{Head of Engineering dimension catalog (50 dimensions).}
\label{tab:engineer_dims}
\small
\begin{tabular}{@{}r l p{7.5cm}@{}}
\toprule
\textbf{\#} & \textbf{Dimension} & \textbf{Description} \\
\midrule
E1  & False Positive Prevention (TNR)     & True negative rate; FPs create developer toil \\
E2  & Precision (PPV)                     & When it says ``vulnerable,'' how often is it right? \\
E3  & Recall (TPR)                        & Fraction of real vulnerabilities detected \\
E4  & F1 Score                            & Balanced precision-recall \\
E5  & MCC (Normalized)                    & MCC normalized to $[0,1]$ \\
E6  & File-Level Accuracy                 & Fraction of TPs with location field populated \\
E7  & Location Accuracy (IoU${>}$0.5)     & Line-level accuracy; devs can jump to correct code \\
E8  & Location Precision                  & $1 -$ over-scope rate; focused code highlights \\
E9  & Actionable Finding Rate             & Fraction of TPs with verdict + CWE + location \\
E10 & CWE Accuracy                        & CWE match; devs need to know the fix type \\
E11 & Wall Time Per Task (avg)            & Average seconds; CI/CD latency \\
E12 & P95 Wall Time                       & 95th percentile latency; worst-case CI delay \\
E13 & SLA Compliance (30s)                & Fraction completing within 30 seconds \\
E14 & SLA Compliance (60s)                & Fraction completing within 60 seconds \\
E15 & Throughput (tasks/min)              & CI pipeline throughput \\
E16 & Cost Per Task                       & Average USD per task \\
E17 & Cost Per PR (projected)             & Cost ${\times}$ 5; projected per pull request \\
E18 & Monthly Cost (projected)            & Cost ${\times}$ 5000; projected monthly at scale \\
E19 & Tokens Per Task                     & Average token consumption \\
E20 & Cost Per Actionable Finding         & USD per fully actionable (3/3) finding \\
E21 & Parse Success Rate                  & FULL parse fraction; clean integration output \\
E22 & Format Compliance                   & FULL / (FULL + PARTIAL); instruction following \\
E23 & Error Rate (inverted)               & $1 -$ error rate; production reliability \\
E24 & Parse Failure Prevention            & $1 -$ parse failure rate \\
E25 & Overall Reliability                 & No errors AND parseable output \\
E26 & Reasoning for Developers            & Fraction with reasoning; devs need explanations \\
E27 & Evidence for Fix (source+sink)      & TPs with source + sink; what to fix \\
E28 & Patch Verification                  & Post-patch accuracy; verifies fixes in PR review \\
E29 & SAST Noise Reduction                & SAST FP rejection; filters existing SAST noise \\
E30 & Developer Trust Score               & PPV ${\times}$ parse success; trustable and parseable \\
E31 & Python Accuracy                     & Language-specific accuracy \\
E32 & JavaScript Accuracy                 & Language-specific accuracy \\
E33 & Java Accuracy                       & Language-specific accuracy \\
E34 & Go Accuracy                         & Language-specific accuracy \\
E35 & Language Parity                     & $1 -$ language accuracy gap \\
E36 & Tool Call Efficiency                & Fraction with ${\leq}5$ tool calls \\
E37 & Turn Efficiency                     & Fraction completing in ${\leq}3$ turns \\
E38 & Wasted Tool Prevention              & $1 -$ (tool calls with zero score) rate \\
E39 & Cost Predictability                 & $1 - CV$ of cost; predictable billing \\
E40 & Worst Language Accuracy             & Min accuracy across languages \\
E41 & Injection Detection                 & Detection rate on injection class \\
E42 & XSS Detection                       & Detection rate on XSS class \\
E43 & Auth Detection                      & Detection rate on auth/access control class \\
E44 & Negative Predictive Value           & Among ``not vuln'' predictions, fraction correct \\
E45 & CI Gate Readiness                   & $0.3{\times}$speed + $0.4{\times}$reliable + $0.3{\times}$structured \\
E46 & Noise-to-Signal Ratio               & PPV; fraction of alerts worth investigating \\
E47 & Developer Productivity Index        & Actionable findings per minute \\
E48 & Full Finding Rate (3/3)             & Fraction of TPs with perfect score \\
E49 & Team Adoption Readiness             & $0.3{\times}$low-FP + $0.2{\times}$fast + $0.25{\times}$reliable + $0.25{\times}$actionable \\
E50 & Overall Engineering Score           & Composite: precision, speed, reliability, actionability, location \\
\bottomrule
\end{tabular}
\end{table*}

\subsection{AI-as-Actor Dimensions}
\label{sec:ai_actor_dims}

The AI-as-Actor lens defines 50 dimensions evaluating the model's fitness for fully autonomous security analysis. Table~\ref{tab:ai_actor_dims} presents the complete catalog.

\begin{table*}[t]
\centering
\caption{AI-as-Actor dimension catalog (50 dimensions).}
\label{tab:ai_actor_dims}
\small
\begin{tabular}{@{}r l p{7.5cm}@{}}
\toprule
\textbf{\#} & \textbf{Dimension} & \textbf{Description} \\
\midrule
AI1  & Autonomous Completion Rate          & Fraction without error or parse failure \\
AI2  & Unsupervised Verdict Accuracy       & Overall verdict accuracy without human review \\
AI3  & MCC (Normalized)                    & MCC normalized to $[0,1]$ \\
AI4  & Full Finding Rate (3/3)             & Fraction of TPs with perfect 3/3 score \\
AI5  & Zero-Shot Accuracy                  & Minimal-prompt accuracy proxy \\
AI6  & Multi-Turn Engagement Rate          & Fraction using ${>}1$ turn; investigative behavior \\
AI7  & Multi-Turn Success Rate             & Accuracy on multi-turn tasks \\
AI8  & Single-Turn Success Rate            & Accuracy on single-turn tasks; quick judgment \\
AI9  & Tool Adoption Rate                  & Fraction of tasks using tools; knows when to investigate \\
AI10 & Tool-Use Effectiveness              & Among tool-using tasks, fraction scoring ${>}0$ \\
AI11 & Tool Efficiency (${\leq}5$ calls)   & Fraction with ${\leq}5$ tool calls \\
AI12 & Tool Call Productivity              & Score earned per tool call \\
AI13 & Navigation Speed                    & $1 - \text{median turns} / 20$; convergence speed \\
AI14 & Wasted Effort Prevention            & $1 -$ (tool calls with zero score) rate \\
AI15 & Structured Output Compliance        & FULL parse rate \\
AI16 & Format Self-Correction              & FULL / non-FAILED; recovers to parseable output \\
AI17 & Error Recovery                      & $1 -$ error rate \\
AI18 & Parse Failure Prevention            & $1 -$ parse failure rate \\
AI19 & Evidence Generation                 & Fraction with evidence chain (source identified) \\
AI20 & Reasoning Generation                & Fraction with reasoning field \\
AI21 & Confidence Calibration              & FULL parse AND correct verdict rate \\
AI22 & Overconfident Miss Prevention       & $1 -$ FN rate among FULL-parsed positive tasks \\
AI23 & CWE Identification                  & CWE accuracy among TPs \\
AI24 & Location Identification             & Location accuracy among TPs \\
AI25 & Complete Analysis Rate              & TPs with verdict + CWE + location \\
AI26 & Recall (TPR)                        & True positive rate \\
AI27 & Specificity (TNR)                   & True negative rate \\
AI28 & Precision (PPV)                     & Positive predictive value \\
AI29 & F1 Score                            & Harmonic mean of precision and recall \\
AI30 & Graceful Degradation                & $1 - |\text{common\_acc} - \text{rare\_acc}|$ \\
AI31 & Cross-Language Autonomy             & $1 -$ language accuracy gap \\
AI32 & Cross-Category Autonomy             & Fraction of categories with ${>}50\%$ accuracy \\
AI33 & Worst Category Floor                & Minimum accuracy across categories \\
AI34 & Critical Vuln Autonomous Detection  & Critical CWE detection without guidance \\
AI35 & Post-Patch Discrimination           & Post-patch task accuracy \\
AI36 & SAST FP Discrimination              & SAST FP rejection; independent judgment \\
AI37 & Cost Efficiency                     & Average cost per task (raw USD) \\
AI38 & Token Efficiency                    & Correct verdicts per 1K tokens \\
AI39 & Wall Time Autonomy ($< 60$s)        & Fraction completing within 60 seconds \\
AI40 & Self-Monitoring Proxy               & FULL parse ${\times}$ correct verdict \\
AI41 & Decision Consistency                & Placeholder (0.5); requires multi-seed runs \\
AI42 & Adversarial Robustness              & Placeholder (0.5); requires adversarial prompts \\
AI43 & Task Type Coverage                  & Average accuracy across all task types \\
AI44 & Autonomous Triage Quality           & $0.4{\times}$verdict + $0.3{\times}$evidence + $0.3{\times}$CWE \\
AI45 & Autonomous Remediation Guidance     & $0.4{\times}$location + $0.3{\times}$evidence + $0.3{\times}$reasoning \\
AI46 & Pipeline Integration Readiness      & $0.4{\times}$parse + $0.3{\times}$error + $0.3{\times}$format \\
AI47 & Human Replacement Index             & $0.4{\times}$F1 + $0.3{\times}$complete + $0.3{\times}$reliable \\
AI48 & Agent Loop Efficiency               & Accuracy / (turns / 5); efficient agent loops \\
AI49 & Autonomous Scalability              & $0.6{\times}$completion + $0.4{\times}$cost\_score \\
AI50 & Overall Autonomy Score              & Composite: completion, quality, tool mastery, analysis, reliability \\
\bottomrule
\end{tabular}
\end{table*}

% ============================================================
% 5. EVALUATION METHODOLOGY
% ============================================================
\section{Evaluation Methodology}
\label{sec:methodology}

\subsection{SecLens Integration}
\label{sec:integration}

The \framework framework is implemented as an extension to \seclens~\citep{haldar2024seclens}, a benchmark for evaluating LLM vulnerability detection capability against confirmed CVEs across multiple CWE categories and programming languages.

\seclens employs a two-layer evaluation architecture:

\begin{itemize}[noitemsep,topsep=2pt]
    \item \textbf{Layer 1 (Code-in-Prompt):} The vulnerable function is provided directly in the prompt. This tests the model's pure reasoning ability without tool use. The model must identify whether the code is vulnerable, classify the CWE, and localize the vulnerability---all from a single code snippet.
    \item \textbf{Layer 2 (Tool-Use):} The model is given access to a sandboxed repository clone and three tools: \texttt{read\_file}, \texttt{search}, and \texttt{list\_dir}. This tests real-world auditing ability, requiring the model to navigate a codebase, identify relevant files, and build a multi-step analysis.
\end{itemize}

Each task in \seclens awards up to 3 points for positive (vulnerable) tasks: 1 point for correct verdict, 1 for correct CWE identification, and 1 for correct location (file path + line range with IoU above threshold). Negative tasks (post-patch code, SAST false positives) award 1 point for correct verdict.

The \framework role-specific dimensions consume the per-task result records produced by \seclens, including:
\begin{itemize}[noitemsep,topsep=2pt]
    \item \textbf{Scores:} verdict (0/1), CWE match (0/1), location match (0/1), total earned
    \item \textbf{Parse result:} status (FULL/PARTIAL/FAILED), parsed output fields (verdict, CWE, location, reasoning, evidence chain with source/sink/flow)
    \item \textbf{Metrics:} cost\_usd, total\_tokens, wall\_time\_s, tool\_calls, turns
    \item \textbf{Task metadata:} task\_type (TRUE\_POSITIVE, POST\_PATCH, SAST\_FALSE\_POSITIVE), task\_category (CWE class), task\_language
\end{itemize}

\subsection{Scoring Pipeline}
\label{sec:pipeline}

The end-to-end scoring pipeline proceeds as follows:

\begin{enumerate}[noitemsep,topsep=2pt]
    \item \textbf{Run evaluation.} Execute \seclens with a specified model, dataset, layer, and prompt preset. This produces a JSONL file of per-task \texttt{TaskResult} records.
    \item \textbf{Compute dimensions.} For each role, compute all 50 dimensions from the result records. Each dimension function takes the full list of results and returns a value.
    \item \textbf{Normalize.} Apply normalization to map all dimension values to $[0, 1]$. Most dimensions are already in this range; unbounded dimensions use cohort-relative min-max normalization.
    \item \textbf{Weight and aggregate.} Multiply each normalized dimension score by its role-specific weight and sum to produce the composite Decision Score (Equation~\ref{eq:decision_score}).
    \item \textbf{Generate report.} Produce a per-role report with the Decision Score, per-category subtotals, and individual dimension scores with visual indicators.
\end{enumerate}

\subsection{Normalization and Aggregation}
\label{sec:normalization}

We distinguish three normalization strategies:

\textbf{Rate dimensions} (e.g., TPR, FPR, parse success rate) are naturally in $[0, 1]$ and require no normalization.

\textbf{Inverted rate dimensions} (e.g., error rate, false negative rate) are computed as $1 - \text{rate}$ so that higher values are always better.

\textbf{Unbounded dimensions} (e.g., cost per task, tokens per task, wall time) are normalized relative to the evaluation cohort using Equation~\ref{eq:normalization}. For dimensions where lower is better (cost, time), the normalization includes inversion. When evaluating a single model in isolation, we use reference ranges derived from the dataset characteristics (e.g., cost capped at \$0.50 per task, wall time capped at 120 seconds).

\textbf{Composite dimensions} (e.g., Overall CISO Confidence, Strategic Value Score) combine multiple sub-dimensions with fixed internal weights. These are already designed to produce values in $[0, 1]$.

\textbf{Placeholder dimensions} (e.g., Cross-Run Consistency, Adversarial Robustness) return a neutral value of 0.5 and are flagged in reports. Their weights are distributed such that the framework remains informative even before multi-run data is available.

% ============================================================
% 6. EXPERIMENTAL DESIGN
% ============================================================
\section{Experimental Design}
\label{sec:experiments}

\subsection{Models Under Evaluation}
\label{sec:models}

We plan to evaluate the following models across both layers and multiple prompt presets:

\textbf{Frontier closed-source models:}
\begin{itemize}[noitemsep,topsep=2pt]
    \item Claude Sonnet 4 and Claude Opus 4 (Anthropic)
    \item GPT-4.1 and GPT-4o (OpenAI)
    \item Gemini 2.5 Pro and Gemini 2.5 Flash (Google)
\end{itemize}

\textbf{Open-source and local models:}
\begin{itemize}[noitemsep,topsep=2pt]
    \item Qwen 3 (various sizes: 8B, 32B, 72B)
    \item Llama 3.1 (70B, 405B)
    \item DeepSeek-V3
    \item CodeLlama 70B
\end{itemize}

This selection spans the capability spectrum from frontier models with tool-use support to smaller open-source models that run locally, enabling analysis of how model scale and architecture affect role-specific evaluations.

\subsection{Dataset}
\label{sec:dataset}

We use the \seclens dataset hosted on HuggingFace (\texttt{sidds020/SecLens}), which contains tasks derived from confirmed CVEs across multiple CWE categories and programming languages. The dataset includes three task types:

\begin{itemize}[noitemsep,topsep=2pt]
    \item \textbf{True Positive (TP):} Pre-patch code containing the confirmed vulnerability.
    \item \textbf{Post-Patch:} The same function after the official fix has been applied.
    \item \textbf{SAST False Positive:} Code flagged by static analysis tools but confirmed as non-vulnerable upon manual review.
\end{itemize}

Table~\ref{tab:dataset_stats} summarizes the planned dataset statistics.

\begin{table}[t]
\centering
\caption{Planned dataset statistics.}
\label{tab:dataset_stats}
\small
\begin{tabular}{@{}l r@{}}
\toprule
\textbf{Attribute} & \textbf{Count} \\
\midrule
Total tasks                    & \placeholder \\
True Positive tasks            & \placeholder \\
Post-Patch tasks               & \placeholder \\
SAST False Positive tasks      & \placeholder \\
Distinct CWE categories        & \placeholder \\
Programming languages          & \placeholder \\
Distinct CVEs                  & \placeholder \\
\bottomrule
\end{tabular}
\end{table}

\subsection{Evaluation Protocol}
\label{sec:protocol}

Each model is evaluated under the following conditions:

\begin{itemize}[noitemsep,topsep=2pt]
    \item \textbf{Both layers:} Layer 1 (code-in-prompt) and Layer 2 (tool-use) when the model supports tool calling.
    \item \textbf{Three prompt presets:} \texttt{base} (standard instructions), \texttt{minimal} (minimal guidance), and \texttt{security\_expert} (detailed security analysis instructions).
    \item \textbf{Parallel evaluation:} Up to 10 concurrent workers per model to measure throughput under realistic conditions.
    \item \textbf{Budget controls:} Maximum cost cap of \$50 per full evaluation run.
\end{itemize}

For each model $\times$ layer $\times$ preset combination, we compute all 250 dimensions and the five role-specific Decision Scores.

% ============================================================
% 7. EXPECTED ANALYSIS
% ============================================================
\section{Expected Analysis}
\label{sec:analysis}

\subsection{Per-Role Decision Scores}
\label{sec:per_role}

The primary output of our framework is a matrix of Decision Scores: models $\times$ roles. Table~\ref{tab:decision_scores} presents the placeholder format for this analysis.

\begin{table*}[t]
\centering
\caption{Per-role Decision Scores (0--100) for evaluated models. Higher is better. Placeholder values shown.}
\label{tab:decision_scores}
\small
\begin{tabular}{@{}l c c c c c c@{}}
\toprule
\textbf{Model} & \textbf{CISO} & \textbf{CAIO} & \textbf{Researcher} & \textbf{Head Eng.} & \textbf{AI Actor} & \textbf{Avg.} \\
\midrule
Claude Sonnet 4          & \placeholder & \placeholder & \placeholder & \placeholder & \placeholder & \placeholder \\
Claude Opus 4            & \placeholder & \placeholder & \placeholder & \placeholder & \placeholder & \placeholder \\
GPT-4.1                  & \placeholder & \placeholder & \placeholder & \placeholder & \placeholder & \placeholder \\
GPT-4o                   & \placeholder & \placeholder & \placeholder & \placeholder & \placeholder & \placeholder \\
Gemini 2.5 Pro           & \placeholder & \placeholder & \placeholder & \placeholder & \placeholder & \placeholder \\
Gemini 2.5 Flash         & \placeholder & \placeholder & \placeholder & \placeholder & \placeholder & \placeholder \\
Qwen 3 72B               & \placeholder & \placeholder & \placeholder & \placeholder & \placeholder & \placeholder \\
Qwen 3 32B               & \placeholder & \placeholder & \placeholder & \placeholder & \placeholder & \placeholder \\
Llama 3.1 70B            & \placeholder & \placeholder & \placeholder & \placeholder & \placeholder & \placeholder \\
DeepSeek-V3              & \placeholder & \placeholder & \placeholder & \placeholder & \placeholder & \placeholder \\
\bottomrule
\end{tabular}
\end{table*}

We hypothesize that role-specific scores will diverge meaningfully: a model ranked first for the CISO may not rank first for the Head of Engineering. This divergence is the core value proposition of our framework---it surfaces information that a single aggregate score cannot.

\subsection{Cross-Role Divergence Analysis}
\label{sec:divergence}

To quantify cross-role disagreement, we define the \emph{Role Divergence Index} (RDI) for a model $m$:

\begin{equation}
\text{RDI}(m) = \max_{r, r'} |D^{(r)}(m) - D^{(r')}(m)|
\label{eq:rdi}
\end{equation}

A high RDI indicates that the model's value depends heavily on the stakeholder perspective. We will analyze:

\begin{itemize}[noitemsep,topsep=2pt]
    \item \textbf{Which model pairs show the highest cross-role rank inversions?} For example, Model A ranks above Model B for the CISO but below for the Engineer.
    \item \textbf{Which role pairs show the most divergence?} We expect CISO vs.\ Engineer and Researcher vs.\ CAIO to show the greatest disagreement.
    \item \textbf{What dimension categories drive divergence?} Risk Profile dimensions are heavily weighted by the CISO but not by the Engineer; Tool-Use dimensions are heavily weighted by the AI-Actor but not by the CISO.
\end{itemize}

Table~\ref{tab:divergence} presents the placeholder format.

\begin{table}[t]
\centering
\caption{Cross-role divergence: pairwise Decision Score correlation and mean rank displacement. Placeholder values.}
\label{tab:divergence}
\small
\begin{tabular}{@{}l c c@{}}
\toprule
\textbf{Role Pair} & \textbf{Pearson $\rho$} & \textbf{Mean $\Delta$Rank} \\
\midrule
CISO -- CAIO             & \placeholder & \placeholder \\
CISO -- Researcher       & \placeholder & \placeholder \\
CISO -- Engineer         & \placeholder & \placeholder \\
CISO -- AI Actor         & \placeholder & \placeholder \\
CAIO -- Researcher       & \placeholder & \placeholder \\
CAIO -- Engineer         & \placeholder & \placeholder \\
CAIO -- AI Actor         & \placeholder & \placeholder \\
Researcher -- Engineer   & \placeholder & \placeholder \\
Researcher -- AI Actor   & \placeholder & \placeholder \\
Engineer -- AI Actor     & \placeholder & \placeholder \\
\bottomrule
\end{tabular}
\end{table}

\subsection{Dimension Sensitivity Analysis}
\label{sec:sensitivity}

We will conduct a sensitivity analysis to understand which individual dimensions most influence the Decision Score for each role. For each dimension $i$ and role $r$, we compute:

\begin{equation}
\text{Impact}_i^{(r)} = w_i^{(r)} \cdot \text{Var}_{m}[s_i^{(m)}]
\label{eq:impact}
\end{equation}

This product of weight and cross-model variance identifies dimensions that both matter to the role (high weight) and differentiate between models (high variance). Dimensions with high impact scores are the key drivers of model selection for that role.

\begin{figure}[t]
\centering
\fbox{\parbox{0.9\linewidth}{\centering\vspace{1.5cm}\textbf{Figure Placeholder}\\\vspace{0.3cm}\small Radar chart showing category-level weight distribution for each of the five roles. Each axis represents one of the ten measurement categories. The five role profiles form distinct shapes, illustrating their different evaluation priorities.\vspace{1.5cm}}}
\caption{Category-level weight distribution radar chart for the five stakeholder roles.}
\label{fig:radar}
\end{figure}

\begin{figure}[t]
\centering
\fbox{\parbox{0.9\linewidth}{\centering\vspace{1.5cm}\textbf{Figure Placeholder}\\\vspace{0.3cm}\small Grouped bar chart comparing Decision Scores across models and roles. Each model has five bars (one per role), illustrating how the same model receives different evaluations from different stakeholder perspectives.\vspace{1.5cm}}}
\caption{Decision Score comparison across models and roles.}
\label{fig:scores}
\end{figure}

% ============================================================
% 8. DISCUSSION
% ============================================================
\section{Discussion}
\label{sec:discussion}

\subsection{Limitations}

\textbf{Weight subjectivity.} The weight vectors, while informed by domain expertise and industry frameworks, involve subjective judgment. Different organizations may prioritize different aspects within the same role. We address this by making the weight profiles modular and customizable---organizations can adjust weights to match their specific context.

\textbf{Placeholder dimensions.} Several dimensions (cross-run consistency, adversarial robustness, prompt sensitivity, reproducibility) require multi-run or adversarial evaluation data that is not available from a single benchmark run. These dimensions return a neutral placeholder value of 0.5. While their weights are modest, filling in these dimensions with real data is a priority for future work.

\textbf{Normalization artifacts.} Cohort-relative normalization means that absolute scores are not comparable across different evaluation cohorts. A model's Decision Score may change when the evaluation cohort changes, even if its raw performance is identical. We provide absolute reference ranges as a mitigation.

\textbf{Single-run evaluation.} Without multi-run data, we cannot estimate confidence intervals on dimension scores. Some dimensions (particularly those computed from small subsets like rare CWEs) may have high variance that is not captured in a single evaluation.

\textbf{Approximations.} Several dimensions use proxies rather than direct measurements. For example, \emph{Mean IoU Estimate} approximates IoU from binary location scores, and \emph{Zero-Guidance Capability} uses overall accuracy as a proxy for minimal-prompt performance.

\subsection{Future Work}

\textbf{Confidence calibration field.} Adding a model-reported confidence field to the \seclens output schema would enable direct confidence calibration analysis, replacing several proxy dimensions with more precise measurements.

\textbf{Multi-run evaluation.} Evaluating each model multiple times with different random seeds would enable computation of cross-run stability, confidence intervals, and reproducibility dimensions.

\textbf{Adversarial testing.} Constructing adversarial prompt variants (obfuscated code, misleading comments, planted red herrings) would fill the adversarial robustness dimensions.

\textbf{Custom role profiles.} Organizations could define custom roles with their own 50-dimension catalogs and weight vectors, extending the framework beyond our five predefined roles. A profile editor tool would facilitate this customization.

\textbf{Temporal analysis.} Tracking dimension scores across model versions (e.g., GPT-4 to GPT-4.1) would enable regression risk assessment and progress tracking.

\textbf{Cross-benchmark integration.} Extending the role-specific evaluation methodology to other security benchmarks (CyberSecEval, SecEval) would enable broader model assessment.

% ============================================================
% 9. CONCLUSION
% ============================================================
\section{Conclusion}
\label{sec:conclusion}

We have presented \framework, a 250-dimension benchmark framework that provides role-specific evaluation of large language models for security vulnerability detection. By defining five stakeholder roles---CISO, CAIO/Head of AI, Security Researcher, Head of Engineering, and AI as Actor---each with 50 tailored dimensions and independently calibrated weight vectors, our framework transforms the same benchmark data into five distinct, actionable evaluations.

The core insight is that model selection for security vulnerability detection is not a one-dimensional optimization problem. A model that maximizes overall accuracy may fail on critical vulnerability classes (unacceptable for the CISO), generate excessive false positives (unacceptable for engineering), or lack the autonomous operation capabilities needed for AI agent deployment. Our framework surfaces these distinctions by construction.

The framework integrates with the \seclens two-layer evaluation pipeline, consuming per-task result records to compute all 250 dimensions without requiring any modifications to the underlying benchmark. Each role's 50 dimensions span ten measurement categories---detection accuracy, vulnerability knowledge depth, localization precision, reasoning quality, operational efficiency, tool-use proficiency, robustness, risk profile, scalability, and autonomy readiness---and aggregate to a composite Decision Score between 0 and 100.

We have presented the complete dimension catalog with computational definitions, the weight assignment methodology, and the composite scoring formulation. Planned experiments across frontier and open-source LLMs will validate the framework's ability to surface meaningful cross-role divergences in model evaluation. We release the full implementation as open-source software to enable organizational adoption and community extension.

% ============================================================
% ACKNOWLEDGMENTS
% ============================================================
\section*{Acknowledgments}

The authors thank the open-source LLM and security research communities for the foundational work that enables this research.

% ============================================================
% REFERENCES
% ============================================================
\bibliographystyle{plainnat}

\begin{thebibliography}{20}

\bibitem[Achiam et~al.(2023)]{achiam2023gpt4}
J.~Achiam, S.~Adler, S.~Agarwal, L.~Ahmad, I.~Akkaya, F.~L. Aleman, D.~Almeida, J.~Altenschmidt, S.~Altman, S.~Anadkat, et~al.
\newblock GPT-4 technical report.
\newblock \emph{arXiv preprint arXiv:2303.08774}, 2023.

\bibitem[Amland(1999)]{amland1999riskbased}
S.~Amland.
\newblock Risk-based testing: Risk analysis fundamentals and metrics for software testing including a financial application case study.
\newblock \emph{Journal of Systems and Software}, 49(2-3):287--295, 1999.

\bibitem[Anthropic(2024)]{anthropic2024claude}
Anthropic.
\newblock The Claude model family.
\newblock Technical report, Anthropic, 2024.

\bibitem[Basili et~al.(1994)]{basili1994gqm}
V.~R. Basili, G.~Caldiera, and H.~D. Rombach.
\newblock The Goal Question Metric approach.
\newblock \emph{Encyclopedia of Software Engineering}, pages 528--532, 1994.

\bibitem[Bhatt et~al.(2023)]{bhatt2023cyberseceval}
M.~Bhatt, S.~Chennabasappa, C.~Nikolaidis, S.~Wan, I.~Evtimov, D.~Gabi, D.~Song, F.~Ahmad, C.~Aber, A.~Terzis, et~al.
\newblock Purple Llama CyberSecEval: A secure coding benchmark for language models.
\newblock \emph{arXiv preprint arXiv:2312.04724}, 2023.

\bibitem[Bhatt et~al.(2024)]{bhatt2024purplellama}
M.~Bhatt, S.~Chennabasappa, Y.~Li, C.~Nikolaidis, D.~Song, S.~Wan, F.~Ahmad, C.~Aschermann, Y.~Chen, D.~Deleo, et~al.
\newblock CyberSecEval 2: A wide-ranging cybersecurity evaluation suite for large language models.
\newblock \emph{arXiv preprint arXiv:2404.13161}, 2024.

\bibitem[Fang et~al.(2024)]{fang2024llmagents}
R.~Fang, R.~Bindu, A.~Gupta, Q.~Zhan, and D.~Kang.
\newblock LLM agents can autonomously hack websites.
\newblock \emph{arXiv preprint arXiv:2402.06664}, 2024.

\bibitem[Gebru et~al.(2021)]{gebru2021datasheets}
T.~Gebru, J.~Morgenstern, B.~Vecchione, J.~W. Vaughan, H.~Wallach, H.~Daum{\'e}~III, and K.~Crawford.
\newblock Datasheets for datasets.
\newblock \emph{Communications of the ACM}, 64(12):86--92, 2021.

\bibitem[{Haldar et~al.}(2024)]{haldar2024seclens}
S.~Haldar, K.~K.~Shrish, and T.~M.
\newblock SecLens: A benchmark for evaluating LLM vulnerability detection.
\newblock Technical report, 2024.

\bibitem[{ISO/IEC}(2011)]{iso25010}
{ISO/IEC}.
\newblock {ISO/IEC 25010:2011} --- Systems and software engineering --- Systems and software Quality Requirements and Evaluation (SQuaRE) --- System and software quality models.
\newblock International Standard, 2011.

\bibitem[Li et~al.(2023)]{li2023seceval}
Z.~Li, C.~Wang, Z.~Liu, H.~Wang, S.~Li, and C.~Gao.
\newblock SecEval: A comprehensive benchmark for evaluating cybersecurity knowledge of foundation models.
\newblock \emph{arXiv preprint arXiv:2312.10505}, 2023.

\bibitem[Liang et~al.(2022)]{liang2022helm}
P.~Liang, R.~Bommasani, T.~Lee, D.~Tsipras, D.~Soylu, M.~Yasunaga, Y.~Zhang, D.~Narayanan, Y.~Wu, A.~Kumar, et~al.
\newblock Holistic evaluation of language models.
\newblock \emph{arXiv preprint arXiv:2211.09110}, 2022.

\bibitem[Mitchell et~al.(2019)]{mitchell2019modelcards}
M.~Mitchell, S.~Wu, A.~Zaldivar, P.~Barnes, L.~Vasserman, B.~Hutchinson, E.~Spitzer, I.~D. Raji, and T.~Gebru.
\newblock Model cards for model reporting.
\newblock In \emph{Proceedings of the Conference on Fairness, Accountability, and Transparency}, pages 220--229, 2019.

\bibitem[Pearce et~al.(2022)]{pearce2022asleep}
H.~Pearce, B.~Ahmad, B.~Tan, B.~Dolan-Gavitt, and R.~Karri.
\newblock Asleep at the keyboard? Assessing the security of GitHub Copilot's code contributions.
\newblock In \emph{2022 IEEE Symposium on Security and Privacy (SP)}, pages 754--768, 2022.

\bibitem[Saaty(1990)]{saaty1990ahp}
T.~L. Saaty.
\newblock How to make a decision: The analytic hierarchy process.
\newblock \emph{European Journal of Operational Research}, 48(1):9--26, 1990.

\bibitem[Siddiq and Santos(2022)]{siddiq2022securityeval}
M.~L. Siddiq and J.~C.~S. Santos.
\newblock SecurityEval dataset: Mining vulnerability examples to evaluate machine learning-based code generation techniques.
\newblock In \emph{Proceedings of the 1st International Workshop on Mining Software Repositories Applications for Privacy and Security}, pages 29--33, 2022.

\bibitem[Team et~al.(2023)]{team2023gemini}
G.~Team, R.~Anil, S.~Borgeaud, Y.~Wu, J.-B. Alayrac, J.~Yu, R.~Sorber, J.~Schalkwyk, A.~M. Dai, A.~Hauth, et~al.
\newblock Gemini: A family of highly capable multimodal models.
\newblock \emph{arXiv preprint arXiv:2312.11805}, 2023.

\bibitem[Tony et~al.(2023)]{tony2023llmseceval}
C.~Tony, M.~Mutas, N.~T. Diem~Ngo, and M.~Ferrara.
\newblock LLMSecEval: A dataset of natural language prompts for security evaluations.
\newblock In \emph{2023 IEEE/ACM International Conference on Mining Software Repositories (MSR)}, pages 588--592, 2023.

\bibitem[Velasquez and Hester(2013)]{velasquez2013mcda}
M.~Velasquez and P.~T. Hester.
\newblock An analysis of multi-criteria decision making methods.
\newblock \emph{International Journal of Operations Research}, 10(2):56--66, 2013.

\bibitem[Zhang et~al.(2024)]{zhang2024reposim}
J.~Zhang, C.~Liu, K.~Chen, and Z.~Su.
\newblock Repository-level vulnerability detection with LLMs: Challenges and opportunities.
\newblock \emph{arXiv preprint arXiv:2401.16185}, 2024.

\end{thebibliography}

\end{document}
